% !PRINTED 513
\weberpage{513}
%<<CONT
the groups $\pi^{-1}Q\pi$ of all these functions, and hence also in their
greatest common divisor $R$. Let the operations
of $R$ be denoted by $\sigma$. If then $\pi$ and $\pi'$ are two
operations that produce the same permutation among the quantities (1),
then $\pi'\pi^{-1}$ will restore the original arrangement of the $\psi$,
and hence be equal to one of the quantities $\sigma$, or
\[ \pi' = \sigma\pi. \]
We therefore have the result:

\begin{quote}
\gesperrt{The permutations of the subsidiary group $R\pi$, and only
these, produce one and the same permutation
among the quantities (1).}
\end{quote}

The degree of the \gesperrt{Galois} group of the resolvent $\varphi(t) = 0$
is therefore equal to the number of these subsidiary groups, i.e.\ equal to the
quotient $p : r$, or the index of the divisor $R$ of $P$.

If $R$ is the identical group, so that $r = 1$, and consequently $\varphi(t) = 0$
is a total resolvent, then the degree of its group is just as high
as the degree of the group of the original equation, and the two
groups are moreover isomorphic. With regard to the group, then,
nothing is gained. The given equation is equivalent to the
resolvent, however different their degrees may be.

If, on the other hand, $Q$ is a normal divisor, so that $R$ is
identical with $Q$, then $\varphi(t) = 0$ is a partial resolvent and the degree
of its group is equal to the index $j$ of the divisor $Q$ of $P$.
After the adjunction of a root of this resolvent, the group
of the original equation reduces to $Q$, and hence to
degree $q$. A splitting of the group has therefore taken place.

If one wishes to form resolvents of the lowest possible
degree, one has to seek out divisors of the group $P$ of the smallest
possible index, and hence of the highest possible degree; but if
at the same time a reduction of the group is to occur, then the
divisors must be normal divisors.

\begin{center}
\S.\,157.

\gesperrt{Reduction of the Galois Group by Adjunction
of Arbitrary Irrationalities.}
\end{center}

According to the \gesperrt{Galois} conception, the problem of solving
an algebraic equation is replaced by another,
%CONT>>
